\bookchapter{Приложение к главе 23: Введение в GitHub — от хостинга кода до платформы сотрудничества}

\chapterintro{В предыдущей главе мы говорили о «программном мышлении», а чтобы воплотить его в жизнь и совместно работать над проектами, не обойтись без одного имени — GitHub. Будь вы студент компьютерных наук, готовящийся к конкурсу, хотите сохранить свои программные работы или просто интересуетесь, чем занимаются разработчики по всему миру — GitHub — платформа, которую необходимо знать. В этом разделе начнём с базовых концепций и поможем вам составить чёткое представление.}

\booksection{Что такое GitHub?}

Простыми словами, GitHub — это\textbf{платформа хостинга кода и сотрудничества}, основанная на Git. Запущена в 2008 году, в 2018-м приобретена Microsoft, до сих пор остаётся крупнейшей в мире платформой хостинга open-source кода: хранит сотни миллионов репозиториев практически на всех основных языках программирования.

Если использовать аналогию:\textbf{Git — это дневник на вашем компьютере, а GitHub — публичная блог-платформа}— вы можете публиковать свои работы, следить за чужими проектами, участвовать в обсуждениях и сотрудничестве. Это не просто «черновик программиста», скорее — «технологический город», совместно строимый разработчиками всего мира — здесь проекты рождаются, развиваются, обсуждаются и переиспользуются.

Для обычных людей GitHub тоже не за семью горами. Многие непрограммисты используют его для хостинга заметок, электронных книг, личных блогов и даже бэкапа конфигурационных файлов. Если у вас есть хоть малейшая потребность «собрать своё и поделиться» — GitHub стоит узнать.

\booksection{Разница между Git и GitHub (ключевые понятия!)}

Многие при первом контакте путают эти два слова. Вот разбор:

\begin{center}
\begin{tabularx}{\linewidth}{|X|X|X|}
\hline
\textbf{Аспект} & \textbf{Git} & \textbf{GitHub} \tabularnewline
\hline
Суть & Инструмент контроля версий (ПО) & Платформа хостинга кода (сайт) \tabularnewline
Где работает & На локальном компьютере & В облаке (доступ через браузер) \tabularnewline
Главная задача & Записывать каждое изменение файлов (что, кто и когда изменил) & Разместить Git-репозиторий в облаке для хранения, демонстрации и сотрудничества \tabularnewline
Нужна ли сеть & Нет, работает офлайн & Нужно подключение для доступа \tabularnewline
Принадлежность & Open-source ПО, создано Linus Torvalds (отцом Linux) & Продукт Microsoft \tabularnewline
Аналоги & SVN, Mercurial & GitLab, Gitee, Bitbucket \tabularnewline
\hline
\end{tabularx}
\end{center}

Простыми словами:\textbf{Git управляет «как записывать изменения», а GitHub — «где хранить изменения и как сотрудничать»}. Git — инструмент, работает локально; GitHub — платформа, размещает Git-репозиторий в облаке для сотрудничества и демонстрации. Можно использовать Git без GitHub (например, для локального контроля версий), но использовать GitHub без Git нельзя (GitHub основан на Git).

\booksection{Ключевые концепции и функции GitHub}

\booksubsection{Базовые концепции}

Прежде чем начать, познакомьтесь с часто встречающимися терминами — иначе интерфейс GitHub покажется непонятным:

\begin{itemize}
  \item \textbf{Репозиторий (Repository, Repo)}: один проект — один репозиторий, в нём хранятся все файлы кода, история версий и документация. Можно считать «проектной папкой».
  \item \textbf{Коммит (Commit)}: запись одного изменения кода. Каждый Commit содержит сообщение, описывающее «что изменено и почему». Это основа «путешествий во времени» Git/GitHub.
  \item \textbf{Ветка (Branch)}: независимая линия разработки, отходящая от основной. Хотите добавить новую функцию — создайте ветку, разрабатывайте в ней, затем слейте обратно в основную, не затрагивая чужую работу.
  \item \textbf{Клонирование (Clone)}: полное копирование облачного репозитория на локальный компьютер; после этого можно изменять код локально.
  \item \textbf{Push / Pull}: Push — загрузка локального кода на GitHub, Pull — скачивание новейшего кода с GitHub на локальную машину. Два ключевых действия для синхронизации при командной работе.
  \item \textbf{Pull Request (PR)}: когда вы изменили код в ветке и хотите слить её в основную — создаёте PR, чтобы команда могла просмотреть код перед слиянием. Основной механизм сотрудничества на GitHub.
  \item \textbf{Issue}: используется для записи багов, запросов функций, тем для обсуждения. Каждый Issue — как пост в обсуждении, можно назначать ответственного, ставить метки.
  \item \textbf{Fork}: копирование чужого репозитория в свой аккаунт; вы можете свободно изменять свой экземпляр, не затрагивая оригинал. Закончив — отправьте автору PR с вашими улучшениями.
\end{itemize}

\booksubsection{Ключевые функции}

GitHub предоставляет не только хранение кода, но и:

\begin{itemize}
  \item \textbf{Контроль версий}: запись истории каждого изменения кода; если что-то сломалось — всегда можно вернуться к предыдущей рабочей версии.
  \item \textbf{Командная работа}: поддержка параллельной разработки несколькими пользователями; одновременная работа не перезаписывает и не конфликтует.
  \item \textbf{Отслеживание проблем}: фиксация багов и запросов функций через Issue для командного контроля.
  \item \textbf{Управление проектом}: встроенные задачи и канбан-доски (Projects) — аналог лёгкого Trello.
  \item \textbf{Автоматизация}: через\textbf{GitHub Actions}— автосборка, тестирование, упаковка и деплой после каждого коммита; ключевой инструмент современного «CI/CD (непрерывной интеграции/непрерывного деплоя)».
  \item \textbf{Управление документацией}: README описывает проект, Wiki хранит подробную архитектуру и стандарты разработки, Pages может бесплатно хостить статический сайт.
  \item \textbf{Хостинг и демонстрация кода}: у каждого репозитория красивая веб-страница с подсветкой синтаксиса — витрина ваших технических работ.
  \item \textbf{Open-source сообщество}: можно Star (в избранное) проекты, Watch (следить) за обновлениями, Follow (подписаться на) разработчиков, участвовать в глобальном open-source сотрудничестве.
\end{itemize}

\booksection{Регистрация и первый шаг: за 5 минут создайте свой первый репозиторий}

\booksubsection{Регистрация}

\begin{enumerate}
  \item Перейдите на\href{https://github.com}{github.com}, нажмите «Sign up».
  \item Укажите email, задайте пароль и имя пользователя (оно появится в адресе вашего профиля, рекомендуется английское, например\texttt{github.\allowbreak com/\allowbreak yourname}).
  \item Пройдите верификацию email и CAPTCHA.
  \item Выберите бесплатный план (Free) — для личных нужд более чем достаточно.
\end{enumerate}

\booksubsection{Создание первого репозитория}

\begin{enumerate}
  \item После входа нажмите «+» в правом верхнем углу, выберите «New repository».
  \item Заполните имя репозитория (Repository name), например\texttt{my-\allowbreak first-\allowbreak project}.
  \item Напишите краткое описание (Description), чтобы другие понимали, что в этом репозитории.
  \item Выберите Public (публичный) или Private (приватный). Публичный виден всем, приватный — только вам и соавторам. Бесплатный аккаунт поддерживает неограниченное число приватных репозиториев.
  \item Отметьте «Add a README file» — README — «паспорт» проекта; его рекомендуется иметь в каждом репозитории.
  \item Нажмите «Create repository» — ваш первый GitHub-репозиторий создан!
\end{enumerate}

\begin{tcbinfo}{Факт: что такое README}
README.md — файл в формате Markdown, первое, что видят при открытии репозитория. Хороший README обычно содержит: описание проекта, инструкции по установке, примеры использования, лицензию, контакты. Грамотный README — базовая компетенция программиста и первое препятствие для того, чтобы проект начали использовать другие.

\end{tcbinfo}

\booksection{Базовый рабочий процесс: от локального к облачному}

Суть рабочего процесса GitHub — запомнить цикл:\textbf{локальное изменение staging commit push обновление в облаке}. Поясним каждый шаг:

\begin{enumerate}
  \item \textbf{Изменить файлы}: пишете код или редактируете документ в локальном редакторе.
  \item \textbf{git add}: добавляете изменённые файлы в «область подготовки» (staging area) — сообщаете Git: «эти изменения я хочу зафиксировать».
  \item \textbf{git commit}: фиксируете подготовленные изменения как «коммит» с пояснительным сообщением, например «исправлена ошибка перехода после логина». Этот шаг выполняется только локально.
  \item \textbf{git push}: загружаете локальные коммиты на GitHub — теперь их видят другие.
  \item \textbf{git pull}: перед началом работы загружаете новейший код из облака, чтобы локальная версия синхронизировалась с командой.
\end{enumerate}

Для новичков запомните правило:\textbf{«сначала pull, потом редактируй; закончив — add, commit с сообщением, и push»}. Следуя этой привычке, командная работа не будет хаотичной.

\booksection{Командный рабочий процесс: как работает Pull Request}

Сила GitHub — в механизме сотрудничества. Типичный командный процесс:

\begin{enumerate}
  \item \textbf{Создать ветку}: от основной ветки (обычно main или master) создайте новую, например\texttt{feature/\allowbreak login}, и разрабатывайте новую функцию в ней.
  \item \textbf{Разрабатывать и коммитить}: пишете код в новой ветке, делаете несколько коммитов по ходу работы.
  \item \textbf{Push на GitHub}: отправляете ветку в облачный репозиторий.
  \item \textbf{Создать Pull Request}: на сайте GitHub создайте PR с запросом на слияние вашей ветки с основной. Опишите, что изменено и почему.
  \item \textbf{Просмотр кода (Code Review)}: команда изучает ваш код, даёт замечания или предложения. Вы можете доработать и закоммитить — PR обновится автоматически.
  \item \textbf{Слияние}: после утверждения нажмите «Merge» — ваш вклад официально входит в проект!
\end{enumerate}

Этот процесс гарантирует, что\textbf{любой код попадает в основную ветку только после проверки}, существенно снижая вероятность ошибок. Open-source проекты используют тот же механизм для принятия вкладов от незнакомцев со всего мира — вы можете отправить PR в известный open-source проект, и если его примут, вы станете контрибутором.

\booksection{Типичные сценарии использования GitHub}

\booksubsection{Студенты: конкурсы и курсовые проекты}

Студенты компьютерных специальностей, участвующие в ACM, хакатонах, дипломных проектах, почти не могут обойтись без GitHub. Это и платформа командной работы, и витрина проектов — жюри может посмотреть ваш репозиторий, оценить качество кода и процесс сотрудничества. Многие компании при найме просят ссылку на профиль GitHub как подтверждение технических способностей.

\booksubsection{Личное хранилище знаний и хостинг заметок}

Непрограммисты тоже используют GitHub. Многие хостят учебные заметки, технические блоги (с GitHub Pages можно бесплатно создать сайт), электронные книги. Markdown-написание + Git-контроль версий + облачный бэкап — три в одном.

\booksubsection{Участие в open-source проектах}

GitHub — центр мирового open-source движения. Можно Fork понравившийся проект, модифицировать и отправить PR обратно в сообщество; или просто Star и отслеживать обновления. Участие в open-source — отличный путь для роста навыков и репутации.

\booksubsection{Автоматический деплой и CI/CD}

С помощью GitHub Actions можно настроить полный автомат: «push кода автотесты автодеплой на сервер». Многие личные блоги и документационные сайты так и построены — написали статью, push — и сайт обновляется сам.

\booksubsection{Бэкап кода и синхронизация между устройствами}

Используйте репозиторий как облако: при смене компьютера или многоустройственной работе — git pull, и у вас новейший код. Никаких флешек.

\booksection{Советы новичкам}

\begin{itemize}
  \item \textbf{Пишите понятные сообщения коммитов}: не пишите «update» — это ни о чём. Хороший commit message: «Исправлена ошибка перехода на неправильную страницу после логина» — и через полгода будет ясно, что делалось.
  \item \textbf{Частые коммиты, маленькими шагами}: закончили мини-функцию — коммит, а не после целого дня работы. Так проще локализовать проблемы.
  \item \textbf{Разрабатывайте в ветках, не правьте основную напрямую}: даже в личном проекте привыкайте к веткам — при командной работе не придётся лихорадочно переучиваться.
  \item \textbf{Используйте .gitignore}: этот файл говорит Git, какие файлы не включать в контроль версий (например, конфиги с паролями, артефакты сборки, папки зависимостей) — чтобы не загружать мусор.
  \item \textbf{Изучайте отличные open-source проекты}: полистайте GitHub Trending, посмотрите структуру кода, README, процесс сотрудничества у высокозвездных проектов — это самый быстрый бесплатный способ обучения.
  \item \textbf{Оформите GitHub Student Developer Pack}: после верификации студенческого статуса бесплатно получаете массу платных инструментов и сервисов (Copilot, домены, облако и др.) на сумму тысяч долларов.
\end{itemize}

\begin{tcbwarn}{Предупреждение безопасности}
\begin{itemize}
  \item Никогда не коммитьте пароли, API-ключи, строки подключения к БД и другую чувствительную информацию в GitHub, особенно в публичные репозитории — утечка будет мгновенно найдена краулерами, последствия серьёзны.
  \item Если чувствительная информация всё-таки попала в коммит, простое удаление файла недостаточно (в истории-то осталось) — используйте\texttt{git filter-\allowbreak branch}или спец-инструменты для очистки истории и немедленно смените утёкшие ключи.
  \item Приватный репозиторий тоже не абсолютно безопасен — чувствительную информацию лучше хранить в переменных окружения, не в коде.
\end{itemize}

\end{tcbwarn}

\begin{tcbinfo}{Доступ к GitHub из Китая}
Доступ к GitHub из Китая иногда нестабилен: страница может не открываться, клонирование — идти медленно. Попробуйте: (1) используйте китайские зеркала (GitClone, ghproxy) для ускорения клонирования; (2) настройте Hosts или используйте прокси; (3) если основная аудитория в Китае, рассмотрите китайскую альтернативу\textbf{Gitee}— функционал аналогичен, скорость доступа выше. Многие команды одновременно ведут репозитории на GitHub и Gitee, совмещая интернациональность и скорость в Китае.

\end{tcbinfo}

\textbf{Продвинутый приём: бесплатный личный блог на GitHub Pages}\par

GitHub Pages позволяет бесплатно хостить статический сайт; в связке с генераторами вроде Jekyll, Hexo, Hugo легко создать личный блог. Написали Markdown-статью, push в репозиторий — блог обновляется автоматически, с HTTPS и кастомным доменом. У бесчисленного множества технических блогеров блог так и устроен: нулевая стоимость, без рекламы, с контролем версий.

\booksection{Итог: почему каждый должен знать GitHub}

GitHub давно перестал быть «инструментом только для программистов» — это скорее\textbf{способ работы цифровой эпохи}: контроль версий делает всё отслеживаемым, механизмы сотрудничества делают команду эффективной, дух open-source позволяет знаниям свободно циркулировать. Даже если вы не пишете код, понимание концепций «коммит, ветка, слияние» поможет лучше управлять любым контентом, требующим «многократных правок + совместной работы» — от диссертаций до документов, от заметок до proposals.

Для студентов компьютерных специальностей GitHub — must-have: это ваше портфолио, платформа сотрудничества, библиотека учебных ресурсов и самое убедительное «резюме» при поиске работы. Регистрируйтесь пораньше, загружайте курсовые и мини-проекты, накапливайте — через несколько лет у вас будет техническая биография, от которой загорятся глаза у интервьюера.

\begin{tcbpractice}{Попробуйте: завершите своё первое знакомство с GitHub}
\begin{enumerate}
  \item Зарегистрируйте аккаунт на GitHub, заполните профиль (аватар, кратко о себе, интересующие технологии).
  \item Создайте публичный репозиторий\texttt{hello-\allowbreak github}, отметьте «Add a README file».
  \item В репозитории нажмите кнопку редактирования README.md (иконка карандаша), напишите самоописание на Markdown, закоммитьте изменения.
  \item В строке поиска GitHub введите интересующее ключевое слово (например,\texttt{python-\allowbreak tutorial}), добавьте два проекта в Star.
  \item (Продвинуто) Установите Git на компьютер, склонируйте репозиторий локально, измените README и push обратно — завершите полный цикл «локально облако».
\end{enumerate}

\end{tcbpractice}
